\begin{abstract}
本文件展示沈阳理工大学博士学位论文中文摘要的排版方式。此处应概括研究背景、需要解决的科学或工程问题以及开展研究的意义。摘要应当能够独立阅读，避免将目录、章节安排或未经解释的缩写直接作为摘要内容。

第二段可说明研究对象、基本假设与总体方法，并围绕主要贡献组织文字。对于每项贡献，建议说明已有方法的局限、本文提出的方法以及得到的主要结论。实验结果应来自真实研究，不能沿用模板中的示例数据。

最后可总结研究结论和适用范围。中文摘要与英文摘要应在研究内容、贡献和结论上相互对应。关键词应准确反映论文主题，具体数量和摘要篇幅请按学校最新要求核对。本段及前述文字均为写作提示，提交前应全部替换。
\keywords{学位论文；排版示例；研究方法；学术写作}
\end{abstract}

\begin{enabstract}
This example demonstrates the layout of an English abstract for a doctoral dissertation at Shenyang Ligong University. Replace this paragraph with the research background, the problem being addressed, and the motivation for the study. The abstract should be self-contained and should define abbreviations when they first appear.

Describe the research methods and principal contributions in a clear sequence. For each contribution, explain the limitation being addressed, the proposed approach, and the evidence supporting the conclusion. Report only results obtained in the actual study; the illustrative data in this template are not research findings.

Conclude with the main findings and their scope of applicability. Ensure that the English and Chinese abstracts describe the same work and that technical terms are used consistently. All text on this page is instructional placeholder content and must be replaced before submission.
\enkeywords{Dissertation; Typesetting example; Research methods; Academic writing}
\end{enabstract}
